\section{Introduction}
\label{sec:intro}

Large language model (LLM) serving latency is dominated by attention
\emph{decode}: after the prompt is processed, tokens are generated one at a time,
and each new token attends over the entire growing key--value (KV) cache. Because
a single query row ($N_q{=}1$) reads a large KV cache to produce one output, decode
is the archetypal memory-bound, low-arithmetic-intensity workload, and it is where
production systems spend most of their wall-clock time. Highly tuned kernels such
as FlashInfer~\cite{flashinfer}, FlashMLA~\cite{flashmla}, and the FlashAttention
lineage~\cite{fa1,fa2,fa3,fa4} are deployed on NVIDIA's Blackwell GPUs, but the
open literature contains no prediction-vs-measured roofline characterization of
attention decode on the latest datacenter part (B300 / GB300, \texttt{sm\_103}).
Production write-ups report headline throughput; they do not publish, per kernel,
what the roofline predicted, what was measured, and why the two differ.

\paragraph{The empty cell}
Figure~\ref{fig:landscape} maps the landscape. FlashAttention-4~\cite{fa4} targets
and benchmarks B200 (\texttt{sm\_100}) for BF16 prefill/training; its kernel is
deployed but never characterized on \texttt{sm\_103}. The closest prior open
roofline of MLA-style decode is Grouped-Latent Attention (GLA) on
H100~\cite{gla}; our delta is exactly \texttt{sm\_103} plus KV-quantization.
SnapMLA~\cite{snapmla} studies FP8 MLA decode on Hopper without a Blackwell
roofline; SageAttention3~\cite{sage3} and Attn-QAT~\cite{attnqat} push FP4
attention but for prefill. Microbenchmark-driven analytical
modeling~\cite{roofmethod} establishes why naive roofline fails on modern GPUs
and owns the method; we cite and apply it. The cell for an open,
prediction-vs-measured decode roofline on \texttt{sm\_103}, across attention
types and precisions, is empty. We fill it with \emph{Roofline-Guided Decode}
(RGD), the study this paper presents.

\begin{figure*}[t]
  \centering
  \includegraphics[width=0.86\textwidth]{v12-related-work-landscape.pdf}
  \caption{The related-work landscape (works surveyed in
  Section~\ref{sec:related}: FA4~\cite{fa4}, GLA~\cite{gla}, SnapMLA~\cite{snapmla},
  SageAttention3~\cite{sage3}, FlashInfer/FlashMLA~\cite{flashinfer,flashmla}).
  Existing work covers prefill, Hopper, or ships production Blackwell decode
  without an open roofline; only the RGD row (this work) combines \texttt{sm\_103}
  decode across MHA$\to$GQA$\to$MLA and FP16$\to$FP8$\to$NVFP4 with a
  prediction-vs-measured roofline. GLA is the closest prior (delta:
  \texttt{sm\_103} + KV-quant).}
  \label{fig:landscape}
\end{figure*}

\paragraph{Contributions}
\begin{enumerate}
  \item The first open, prediction-vs-measured roofline characterization of
  attention decode on B300 / \texttt{sm\_103}, spanning three attention types
  (MHA $\to$ GQA $\to$ MLA) and three precisions (FP16 $\to$ FP8 $\to$ NVFP4),
  built as twelve single-variable kernel steps.
  \item A two-layer prediction model (pure roofline plus a per-CTA-corrected
  layer) whose predictions are \emph{pre-registered before} each kernel is coded
  or measured. The gap between the layers is itself the finding: across the seven
  decode steps the per-CTA-corrected layer matched the measured limiter every
  time, while pure roofline predicted the wrong regime at every single-stream
  step (Section~\ref{sec:method}).
  \item The central finding: decode's per-CTA ``wall'' is
  \emph{work-starvation}, not a fixed architectural limit. The right engine
  converts work into throughput. A warp-per-head CUDA-core MLA kernel is flat at
  $0.75\tflops$, whereas the production \texttt{tcgen05} kernel spans
  $2.4 \to 1785\tflops$ scaling with batch $\times$ context. Even so, the roofs
  are reachable only at serving scale, not in common single-stream decode.
\end{enumerate}

\paragraph{Honest scope}
We \emph{complement}, not beat, production kernels: FlashInfer, FlashMLA, and the
CUTLASS reference kernel already run B300 decode and are faster by construction.
The prediction-vs-measured methodology is also not novel; prior work owns
it~\cite{roofmethod}. Our contribution is the open application of that methodology
to this workload $\times$ architecture $\times$ precision space, the
per-CTA-corrected model, and the boundary results (including a negative one,
Section~\ref{sec:stagea}). Figure~\ref{fig:positioning} states precisely what we
claim and what we do not. Prediction misses are reported as first-class results
throughout.
